\chapter{Introduction to Medical Imaging}

\section*{Patient focus}
This chapter explains the examination or clinical problem in accessible language. Patients should use it to prepare questions and understand the purpose, experience, benefits, and limitations of imaging; local instructions and the official report take priority.


\label{ch:intro}

Medical imaging is a cornerstone of modern clinical diagnosis, oncology staging, surgical planning, and therapeutic monitoring. By translating physical energy interactions with human tissues into spatial representations, medical imaging allows clinicians to visualize internal anatomy, physiological processes, and cellular metabolism non-invasively.

---

\section*{1. Classification of Medical Imaging Modalities}
Diagnostic imaging modalities can be classified using three primary criteria:

\subsection*{By Physical Energy Source}
\begin{itemize}
\item \textbf{Electromagnetic Radiation (X-rays)}: High-energy photons are projected through the body. Modalities include planar X-ray radiography, fluoroscopy, mammography, and computed tomography (CT).
\item \textbf{Mechanical Sound Waves}: High-frequency sound waves propagate through soft tissue, reflecting at acoustic boundaries. Modalities include grayscale ultrasound, Doppler flow analysis, and echocardiography.
\item \textbf{Magnetic Fields and Radiofrequency (RF) Waves}: Protons in tissue align with a magnetic field and are excited by RF pulses. Modalities include structural MRI, functional MRI (fMRI), and diffusion-weighted MRI.
\item \textbf{Radioactive Decay (Radiopharmaceuticals)}: Unstable isotopes bound to metabolic ligands are injected and detected. Modalities include SPECT and positron emission tomography (PET).
\end{itemize}

\subsection*{By Ionization Hazard}
\begin{itemize}
\item \textbf{Ionizing Modalities}: Utilize radiation of sufficient energy to liberate electrons from atoms, posing a biological risk of DNA damage. Includes X-ray, fluoroscopy, CT, SPECT, and PET.
\item \textbf{Non-Ionizing Modalities}: Do not utilize ionizing radiation. Includes ultrasound and MRI. Modality choice should be individualized, especially during pregnancy and in children.
\end{itemize}

\subsection*{By Clinical Objective: Structural vs. Functional}
\begin{itemize}
\item \textbf{Structural Imaging}: Focuses on detailed anatomical localization, boundaries, and physical densities (e.g., bones on CT, soft tissue borders on MRI, gallstones on ultrasound).
\item \textbf{Functional / Molecular Imaging}: Focuses on physiological activity, blood perfusion, receptor binding, and cellular metabolism (e.g., glucose consumption on FDG-PET, brain perfusion on fMRI, or ventilation-perfusion mismatch on V/Q scan).
\end{itemize}

---

\section*{2. The Clinical Imaging Workflow}
The imaging pathway is a multi-step clinical process requiring coordination between referring physicians, technologists, medical physicists, and radiologists:

\begin{enumerate}
    \item \textbf{Clinical Indication \& Question}: The referring physician identifies a clinical symptom or diagnostic question (e.g., ``Sudden onset left-sided weakness; rule out stroke'').
    \item \textbf{Modality Selection \& Appropriateness}: Selecting the optimal test. Guidelines (such as the ACR Appropriateness Criteria) help balance diagnostic yield against radiation dose, cost, contrast safety, and speed.
    \item \textbf{Patient Preparation \& Screening}: Screening for contraindications (e.g., checking renal function for contrast safety, or screening for metallic implants before MRI).
    \item \textbf{Image Acquisition}: The radiographer/technologist positions the patient and operates the scanner using optimized protocols.
    \item \textbf{Reconstruction \& Post-Processing}: Fusing datasets, generating reformations (axial, coronal, sagittal), or correcting artifacts.
    \item \textbf{Interpretation}: The radiologist systematically reviews the images, correlates findings with clinical history, and generates a structured report.
    \item \textbf{Clinical Communication}: Communicating findings to the clinical team, including immediate direct verbal reporting of critical or life-threatening results.
\end{enumerate}

\clearpage
